% ============================================================
%  [*] preamble — the look and the vocabulary of the document
%      (Lithuanian edition: the same preamble as docs/latex/en,
%      with the language, the captions, the references and the
%      printed labels in Lithuanian — see the block at the end)
%
%  Everything the chapters share: the page, the typefaces, the
%  colours of the product, the inline macros, the boxes that
%  mark who does what (the model, the code, a person), the
%  listing styles and the drawing styles of the figures.
%
%  The chapters load no package and define no colour of their
%  own: what a chapter needs is defined here, once, so that
%  two chapters cannot draw the same thing two ways.
%
%  Engine: LuaLaTeX (Unicode text, OpenType fonts). Built by
%  docs/buildDocs.sh in the TeX Live container.
% ============================================================


% ------------------------------------------------------------
% Language, type and page
% ------------------------------------------------------------
\usepackage{etoolbox}
\usepackage{fontspec}
\usepackage[shorthands=off,british,lithuanian]{babel}
\usepackage{microtype}
% colon=literal: the inline code macros below are set in math
% mode (the url package), where unicode-math would otherwise
% turn ":" into the ratio sign U+2236, which the code face
% lacks — "gateway:azure/…" would lose its colon.
\usepackage[colon=literal]{unicode-math}

% The product's own faces: the reports it prints are set in
% Source Serif 4 and Source Sans 3 as well.
\setmainfont{SourceSerif4}[
  Extension      = .otf,
  UprightFont    = *-Regular,
  ItalicFont     = *-RegularIt,
  BoldFont       = *-Semibold,
  BoldItalicFont = *-SemiboldIt,
  Numbers        = {Proportional, Lining}]
\setsansfont{SourceSans3}[
  Extension      = .otf,
  UprightFont    = *-Regular,
  ItalicFont     = *-RegularIt,
  BoldFont       = *-Semibold,
  BoldItalicFont = *-SemiboldIt,
  Scale          = MatchLowercase]
\setmonofont{SourceCodePro}[
  Extension      = .otf,
  UprightFont    = *-Regular,
  ItalicFont     = *-RegularIt,
  BoldFont       = *-Semibold,
  BoldItalicFont = *-SemiboldIt,
  Scale          = 0.86,
  HyphenChar     = None]

\usepackage[a4paper, left=30mm, right=30mm, top=27mm, bottom=30mm,
            headheight=14pt, headsep=9mm, footskip=12mm]{geometry}
\usepackage{setspace}
\setstretch{1.08}
\raggedbottom
\emergencystretch=2em
\clubpenalty=8000
\widowpenalty=8000


% ------------------------------------------------------------
% Colours — the product's palette (vite/app/src/brand.js)
% ------------------------------------------------------------
\usepackage[table]{xcolor}
\definecolor{ink}{HTML}{173A6E}        % the brand blue: headings, rules, links
\definecolor{mid}{HTML}{3E6FB0}        % the mark's middle blue
\definecolor{pale}{HTML}{8FB1E3}       % the mark's pale blue
\definecolor{sun}{HTML}{FFD84A}        % the mark's yellow: wherever a model speaks
\definecolor{tint}{HTML}{EEF3FB}       % a blue wash for panels
\definecolor{suntint}{HTML}{FFF7D6}    % a yellow wash for the model's panels
\definecolor{paper}{HTML}{FFFFFF}
\definecolor{body}{HTML}{1B2330}       % running text
\definecolor{muted}{HTML}{5B6676}      % captions, secondary text
\definecolor{rule}{HTML}{D5DCE8}       % hairlines
\definecolor{codebg}{HTML}{F6F8FB}     % behind listings
% The four tones of a verdict, as the console shows them
\definecolor{good}{HTML}{1E7B4E}
\definecolor{goodbg}{HTML}{EDF6F0}
\definecolor{info}{HTML}{1F6FB2}
\definecolor{infobg}{HTML}{EBF2FA}
\definecolor{warn}{HTML}{9A5400}
\definecolor{warnbg}{HTML}{FCF3E5}
\definecolor{bad}{HTML}{B3261E}
\definecolor{badbg}{HTML}{FBEDEB}
\color{body}


% ------------------------------------------------------------
% Headings, running heads, captions, lists
% ------------------------------------------------------------
\setkomafont{disposition}{\sffamily\bfseries\color{ink}}
\setkomafont{chapter}{\sffamily\bfseries\Huge\color{ink}}
\setkomafont{section}{\sffamily\bfseries\Large\color{ink}}
\setkomafont{subsection}{\sffamily\bfseries\large\color{ink}}
\setkomafont{subsubsection}{\sffamily\bfseries\normalsize\color{ink}}
\setkomafont{paragraph}{\sffamily\bfseries\normalsize\color{body}}
\setkomafont{part}{\sffamily\bfseries\Huge\color{ink}}
\setkomafont{partnumber}{\sffamily\bfseries\LARGE\color{mid}}
\setkomafont{descriptionlabel}{\sffamily\bfseries\color{ink}}
\RedeclareSectionCommand[beforeskip=0pt, afterskip=1.6\baselineskip]{chapter}
\RedeclareSectionCommand[beforeskip=1.5\baselineskip, afterskip=.5\baselineskip]{section}
\RedeclareSectionCommand[beforeskip=1.1\baselineskip, afterskip=.3\baselineskip]{subsection}
\RedeclareSectionCommand[beforeskip=.9\baselineskip, afterskip=.2\baselineskip]{subsubsection}
\setcounter{secnumdepth}{2}
\setcounter{tocdepth}{1}

\usepackage[automark, headsepline=.4pt]{scrlayer-scrpage}
\clearpairofpagestyles
\setkomafont{pageheadfoot}{\sffamily\footnotesize\color{muted}}
\setkomafont{pagenumber}{\sffamily\footnotesize\bfseries\color{ink}}
\setkomafont{headsepline}{\color{rule}}
\ihead{\leftmark}
\ohead{Atitikties asistentas}
\ofoot*{\pagemark}
\renewcommand*{\chaptermarkformat}{\thechapter\enskip}

\usepackage{caption}
\captionsetup{font={small,sf}, labelfont={bf,color=ink}, textfont={color=muted},
              labelsep=quad, justification=raggedright, singlelinecheck=false,
              skip=7pt}
\captionsetup[table]{position=top}

\usepackage{enumitem}
\setlist{itemsep=.25\baselineskip, parsep=0pt, topsep=.4\baselineskip}
\setlist[itemize,1]{label={\color{mid}\rule[.22ex]{.42em}{.42em}}, leftmargin=1.4em}
\setlist[itemize,2]{label={\color{pale}\rule[.25ex]{.36em}{.36em}}, leftmargin=1.3em}
\setlist[enumerate,1]{label={\sffamily\bfseries\color{ink}\arabic*.}, leftmargin=1.7em}
\setlist[enumerate,2]{label={\sffamily\color{ink}\alph*)}, leftmargin=1.6em}
\setlist[description]{font=\sffamily\bfseries\color{ink}, leftmargin=1.4em, style=nextline}


% ------------------------------------------------------------
% Tables
% ------------------------------------------------------------
\usepackage{array, booktabs, longtable, tabularx, xltabular, multirow, makecell}
\usepackage{float}
\usepackage{graphicx}
\usepackage{pdflscape}
\renewcommand{\arraystretch}{1.22}
\setlength{\heavyrulewidth}{.9pt}
\setlength{\lightrulewidth}{.4pt}
\arrayrulecolor{rule}
% Column types: L{w} / C{w} / R{w} are ragged paragraphs of a
% set width; Y is tabularx's X without justification.
\newcolumntype{L}[1]{>{\raggedright\arraybackslash}p{#1}}
\newcolumntype{C}[1]{>{\centering\arraybackslash}p{#1}}
\newcolumntype{R}[1]{>{\raggedleft\arraybackslash}p{#1}}
\newcolumntype{Y}{>{\raggedright\arraybackslash}X}
% A header cell: \hd{Text}
\newcommand{\hd}[1]{\textcolor{ink}{\sffamily\bfseries #1}}
% Tables are set a size smaller than the text
\AtBeginEnvironment{tabular}{\small}
\AtBeginEnvironment{tabularx}{\small}
\AtBeginEnvironment{longtable}{\small}
\AtBeginEnvironment{xltabular}{\small}


% ------------------------------------------------------------
% Figures: TikZ and plots
% ------------------------------------------------------------
\usepackage{tikz}
\usetikzlibrary{arrows.meta, positioning, calc, fit, backgrounds, shapes.geometric,
                shapes.misc, shapes.multipart, matrix, chains, decorations.pathreplacing}
\usepackage{pgfplots}
\pgfplotsset{compat=1.18}
\usepgfplotslibrary{groupplots}

% One visual language for every diagram:
%   svc    a service or a module of ours        (white, ink outline)
%   store  a table or a file store              (cylinder)
%   ext    something outside the system         (dashed outline)
%   ai     a model call                         (yellow)
%   chk    a check or a computation by code     (blue wash)
%   man    a decision of a person               (ink, white text)
%   data   a piece of data passed along         (grey, slanted)
%   note   a remark beside the drawing
%   flow   the main path        call   a request and its answer
%   opt    a path taken only sometimes
\tikzset{
  every picture/.style={font=\sffamily\footnotesize, >={Stealth[length=5pt,width=4pt]}},
  box/.style={draw=ink, line width=.6pt, rounded corners=2.5pt, fill=paper, text=body,
              align=center, inner xsep=6pt, inner ysep=5pt, minimum height=8mm,
              execute at begin node={\hyphenpenalty=10000 \exhyphenpenalty=10000}},
  svc/.style={box},
  ext/.style={box, dashed, draw=muted, text=muted},
  ai/.style={box, draw=ink, fill=sun},
  chk/.style={box, fill=tint},
  man/.style={box, fill=ink, text=white, draw=ink,
              execute at begin node={\hypersetup{linkcolor=white}}},
  data/.style={box, draw=muted, fill=codebg, trapezium, trapezium left angle=78,
               trapezium right angle=102, trapezium stretches=true, rounded corners=0pt},
  store/.style={draw=ink, line width=.6pt, fill=paper, cylinder, shape border rotate=90,
                aspect=.22, align=center, inner xsep=6pt, inner ysep=4pt, minimum height=9mm},
  note/.style={text=muted, align=left, font=\sffamily\scriptsize, inner sep=2pt},
  lane/.style={draw=rule, line width=.5pt, fill=tint!55, rounded corners=4pt, inner sep=7pt},
  lanelabel/.style={text=mid, font=\sffamily\scriptsize\bfseries, inner sep=2pt},
  flow/.style={->, draw=ink, line width=.7pt},
  call/.style={->, draw=mid, line width=.6pt, dashed},
  opt/.style={->, draw=muted, line width=.6pt, densely dotted},
  lbl/.style={font=\sffamily\scriptsize, text=muted, fill=paper, inner sep=1.5pt},
}
% The legend of the three actors, for a figure that mixes them:
%   \actorlegend   (place it under the picture, in the figure)
\newcommand{\actorlegend}{%
  \par\smallskip{\sffamily\scriptsize\color{muted}%
  \tikz[baseline=-.6ex]\node[ai, minimum height=3mm, inner sep=2.5pt]{};~modelio iškvietimas\quad
  \tikz[baseline=-.6ex]\node[chk, minimum height=3mm, inner sep=2.5pt]{};~kodas tikrina ar skaičiuoja\quad
  \tikz[baseline=-.6ex]\node[man, minimum height=3mm, inner sep=2.5pt]{};~sprendžia žmogus\quad
  \tikz[baseline=-.6ex]\node[svc, minimum height=3mm, inner sep=2.5pt]{};~žingsnis be modelio}}


% ------------------------------------------------------------
% Inline vocabulary
% ------------------------------------------------------------
% \code{run_prompt}      an identifier: function, constant, field, slug, value
% \file{analysis/pipeline.py}   a path in the repository
% \route{POST}{/api/reports}    an HTTP route
% \prompt{report.gap_analysis}  a prompt by its slug
% \model{gpt-6.1-sol}           a model by the name the gateway gives it
% \ui{Veiksmai}                 a label as the console prints it
% \lt{atitikties pareigūnas}    Lithuanian words inside English text
% \term{finding}                a term where it is defined
% \status{met}                  one of a finding's five statuses, as a chip
% \chip[colour]{text}           any other small tag
\usepackage[obeyspaces,spaces]{url}
\urlstyle{tt}
\Urlmuskip=0mu
% Names of the code are set like URLs: any character may stand
% in them, and a long one breaks after a slash, a dot or an
% underscore instead of running into the margin. A model's
% name is never broken. obeyspaces keeps a space inside the
% argument (by default the url package drops it, so a command
% or an SQL clause would print run together); spaces lets the
% line break there.
\DeclareUrlCommand\codeurl{\urlstyle{tt}}
\DeclareUrlCommand\codeurlwhole{\urlstyle{tt}\def\UrlBreaks{}\def\UrlBigBreaks{}}
\DeclareRobustCommand{\code}{\codeurl}
\DeclareRobustCommand{\file}{\codeurl}
\DeclareRobustCommand{\prompt}{\codeurl}
\DeclareRobustCommand{\model}{\codeurlwhole}
\DeclareRobustCommand{\route}[1]{\textcolor{mid}{\sffamily\footnotesize\bfseries #1}\,\codeurl}
\DeclareRobustCommand{\ui}[1]{\textcolor{ink}{\sffamily #1}}
\DeclareRobustCommand{\lt}[1]{#1}
% An English term at its first use: (angl. \en{prompt injection})
\DeclareRobustCommand{\en}[1]{\foreignlanguage{british}{\emph{#1}}}
\DeclareRobustCommand{\term}[1]{\emph{#1}}

\usepackage[most]{tcolorbox}
\tcbset{boxrule=.5pt, arc=2pt, left=8pt, right=8pt, top=6pt, bottom=6pt}
\newtcbox{\chip}[1][ink]{on line, arc=2pt, colback=#1!9, colframe=#1!45, boxrule=.35pt,
  boxsep=0pt, left=3.2pt, right=3.2pt, top=1.4pt, bottom=1.4pt,
  fontupper=\sffamily\scriptsize\bfseries\color{#1}}
\DeclareRobustCommand{\status}[1]{%
  \ifstrequal{#1}{met}{\chip[good]{Atitinka}}{%
  \ifstrequal{#1}{partial}{\chip[warn]{Iš dalies}}{%
  \ifstrequal{#1}{missing}{\chip[bad]{Neatitinka}}{%
  \ifstrequal{#1}{unknown}{\chip[muted]{Neįrodyta}}{%
  \chip[muted]{Netaikoma}}}}}}


% ------------------------------------------------------------
% Boxes that say who is acting
% ------------------------------------------------------------
% aibox      what a model is asked and what it answers
% codebox    what the code checks or computes, with no model
% humanbox   what only a person decides
% lessonbox  a measurement or a failure, with its date, and what
%            changed because of it
% notebox    anything else worth setting apart
% Each takes its title as the mandatory argument.
\newtcolorbox{aibox}[1]{enhanced, breakable, colback=suntint, colframe=sun!70!warn!60,
  borderline west={2.6pt}{0pt}{sun}, boxrule=0pt, frame hidden, sharp corners,
  left=11pt, title={#1}, coltitle=body, fonttitle=\sffamily\bfseries\small,
  attach title to upper={\par\smallskip}, before skip=10pt, after skip=10pt}
\newtcolorbox{codebox}[1]{enhanced, breakable, colback=tint, colframe=mid,
  borderline west={2.6pt}{0pt}{mid}, boxrule=0pt, frame hidden, sharp corners,
  left=11pt, title={#1}, coltitle=ink, fonttitle=\sffamily\bfseries\small,
  attach title to upper={\par\smallskip}, before skip=10pt, after skip=10pt}
\newtcolorbox{humanbox}[1]{enhanced, breakable, colback=paper, colframe=ink,
  boxrule=.6pt, leftrule=2.6pt, sharp corners,
  left=11pt, title={#1}, coltitle=ink, fonttitle=\sffamily\bfseries\small,
  attach title to upper={\par\smallskip}, before skip=10pt, after skip=10pt}
\newtcolorbox{lessonbox}[1]{enhanced, breakable, colback=codebg, colframe=rule,
  borderline west={2.6pt}{0pt}{pale}, boxrule=0pt, frame hidden, sharp corners,
  left=11pt, title={#1}, coltitle=ink, fonttitle=\sffamily\bfseries\small,
  attach title to upper={\par\smallskip}, before skip=10pt, after skip=10pt}
\newtcolorbox{notebox}[1]{enhanced, breakable, colback=paper, colframe=rule,
  boxrule=.5pt, title={#1}, coltitle=ink, fonttitle=\sffamily\bfseries\small,
  attach title to upper={\par\smallskip}, before skip=10pt, after skip=10pt}

% A breakable box may start so low on a page that only its title
% fits there and its text goes on overleaf. Every box with a title
% (these five, the path card, a titled code file and a prompt
% excerpt) therefore asks for six lines of room before it begins;
% with less left, it starts on the next page. In a table cell the
% request does nothing.
\usepackage{needspace}
\newcommand{\keepboxtitle}{\par\needspace{6\baselineskip}}
\AddToHook{env/aibox/before}{\keepboxtitle}
\AddToHook{env/codebox/before}{\keepboxtitle}
\AddToHook{env/humanbox/before}{\keepboxtitle}
\AddToHook{env/lessonbox/before}{\keepboxtitle}
\AddToHook{env/notebox/before}{\keepboxtitle}
\AddToHook{env/pathcard/before}{\keepboxtitle}
\AddToHook{env/codefile/before}{\keepboxtitle}
\AddToHook{env/promptexcerpt/before}{\keepboxtitle}

% The card of one AI path: the same ten questions for every
% place where a model is called.
%   \begin{pathcard}{Gap analysis}{report.gap_analysis}
%     \pcrow{Trigger}{...}
%     ...
%   \end{pathcard}
% The rows, in this order: Trigger, Input, Prompt and model,
% Output, Checked by code, Stored, On failure, Cost, Human
% control, Measured by.
\newenvironment{pathcard}[2]{%
  \begin{tcolorbox}[enhanced, breakable, colback=paper, colframe=ink, boxrule=.6pt, arc=2pt,
    colbacktitle=ink, coltitle=white, fonttitle=\sffamily\bfseries\small,
    title={#1\hfill{\mdseries\footnotesize\codeurlwhole{#2}}},
    left=8pt, right=8pt, top=4pt, bottom=4pt, before skip=12pt, after skip=12pt]
  \small
  \begin{description}[style=multiline, leftmargin=31mm, labelwidth=29mm, labelsep=2mm,
    font=\sffamily\bfseries\footnotesize\color{ink}, itemsep=3pt, topsep=1pt]}%
  {\end{description}\end{tcolorbox}}
% The ten rows keep their English keys in the source; the card prints
% them in Lithuanian, so that every card in the book reads the same.
\newcommand{\pcrowname}[1]{%
  \ifstrequal{#1}{Trigger}{Paleidimas}{%
  \ifstrequal{#1}{Input}{Įvestis}{%
  \ifstrequal{#1}{Prompt and model}{Instrukcija ir modelis}{%
  \ifstrequal{#1}{Output}{Išvestis}{%
  \ifstrequal{#1}{Checked by code}{Kodo patikros}{%
  \ifstrequal{#1}{Stored}{Saugoma}{%
  \ifstrequal{#1}{On failure}{Nepavykus}{%
  \ifstrequal{#1}{Cost}{Kaina}{%
  \ifstrequal{#1}{Human control}{Žmogaus kontrolė}{%
  \ifstrequal{#1}{Measured by}{Kaip matuojama}{#1}}}}}}}}}}}
\newcommand{\pcrow}[2]{\item[\pcrowname{#1}] #2}


% ------------------------------------------------------------
% Code and prompt text, word for word
% ------------------------------------------------------------
% Set with fancyvrb (fvextra), not listings: listings moves
% every non-ASCII character to the front of its word under a
% Unicode engine, and the prompts and the data are full of
% Lithuanian. Long lines wrap; nothing is highlighted.
%
%   \begin{codeblock}                    a piece of code or data
%   \begin{codefile}{what it shows}      the same under a title line
%   \begin{promptexcerpt}{report.gap_analysis}{what the passage is about}
%                                        a passage of a prompt, word for word
% The \begin line and the \end line stand alone on their lines.
% None of the three takes an optional argument: looking ahead
% for one would read the first line of the text as TeX.
\usepackage{fvextra}
\usepackage[style=german]{csquotes}     % after fvextra, as fvextra asks; „…“ as Lithuanian sets them
\fvset{fontsize=\footnotesize, breaklines=true, breakafter={/,_}, breakaftersymbolpre={},
       breaksymbolleft={\tiny\ensuremath{\hookrightarrow}}, breakindent=1.2em, tabsize=2,
       formatcom=\color{body},
       listparameters={\setlength{\topsep}{0pt}\setlength{\partopsep}{0pt}\setlength{\parskip}{0pt}}}
\tcbset{
  codestyle/.style={enhanced, breakable, colback=codebg, boxrule=0pt, frame hidden,
    sharp corners, borderline west={2pt}{0pt}{rule}, left=9pt, right=6pt, top=3pt, bottom=3pt,
    before skip=9pt, after skip=9pt, coltitle=muted, fonttitle=\sffamily\footnotesize,
    attach title to upper={\par\vspace{1pt}}},
  promptstyle/.style={enhanced, breakable, colback=suntint, boxrule=0pt, frame hidden,
    sharp corners, borderline west={2.6pt}{0pt}{sun}, left=11pt, right=8pt, top=4pt, bottom=4pt,
    before skip=10pt, after skip=10pt, coltitle=body, fonttitle=\sffamily\footnotesize,
    attach title to upper={\par\vspace{2pt}}}}
\NewDocumentEnvironment{codeblock}{}{%
  \VerbatimEnvironment
  \begin{tcolorbox}[codestyle]%
  \begin{Verbatim}}%
  {\end{Verbatim}\end{tcolorbox}}
\NewDocumentEnvironment{codefile}{m}{%
  \VerbatimEnvironment
  \begin{tcolorbox}[codestyle, title={#1}]%
  \begin{Verbatim}}%
  {\end{Verbatim}\end{tcolorbox}}
\NewDocumentEnvironment{promptexcerpt}{m m}{%
  \VerbatimEnvironment
  \begin{tcolorbox}[promptstyle,
    title={{\bfseries Iš instrukcijos} {\codeurlwhole{#1}}\quad{\color{muted}#2}}]%
  \begin{Verbatim}[breaksymbolleft={}, breakindent=0pt, breakafter={}]}%
  {\end{Verbatim}\end{tcolorbox}}


% ------------------------------------------------------------
% Links and references
% ------------------------------------------------------------
\usepackage[colorlinks=true, linkcolor=ink, urlcolor=mid, citecolor=ink,
            bookmarksnumbered=true, pdfusetitle]{hyperref}
\usepackage[nameinlink]{cleveref}
\pdfstringdefDisableCommands{\def\code#1{\detokenize{#1}}\def\file#1{\detokenize{#1}}%
  \def\prompt#1{\detokenize{#1}}\def\model#1{\detokenize{#1}}%
  \def\route#1#2{#1 \detokenize{#2}}\def\ui#1{#1}\def\lt#1{#1}\def\term#1{#1}\def\status#1{#1}}
% (the Lithuanian reference formats are set at the end)


% ------------------------------------------------------------
% Small conveniences
% ------------------------------------------------------------
\newcommand{\product}{Atitikties asistentas}
\newcommand{\productlt}{Atitikties asistentas}
\newcommand{\demo}{1009~LTD}
\newcommand{\eg}{pvz.\ }
\newcommand{\ie}{t.~y.\ }
% The product's mark: four tiles, two of them quarter discs.
% \mosaic{scale}
\newcommand{\mosaic}[1]{%
  \begin{tikzpicture}[x=#1, y=#1, baseline=0]
    \fill[ink] (31,-31) -- (8,-31) arc[start angle=180, end angle=90, radius=23] -- cycle;
    \fill[mid, rounded corners=4*#1] (33,-8) rectangle (56,-31);
    \fill[mid, rounded corners=4*#1] (8,-33) rectangle (31,-56);
    \fill[sun] (33,-33) -- (56,-33) arc[start angle=0, end angle=-90, radius=23] -- cycle;
  \end{tikzpicture}}



% ------------------------------------------------------------
% Lithuanian captions, parts and references
% ------------------------------------------------------------
% A caption reads "3.2 pav. …" and "2.1 lentelė …", a part
% "I dalis". A reference puts the number first and an
% abbreviation after it, which reads right in every case of
% the noun: (9 sk.), žr. 3.2 pav., 2.1 lent., 9.3 posk.
% An appendix is named in full, "A priedas": refer to one in
% parentheses, (\cref{app:prompts}).
\renewcommand*{\figureformat}{\thefigure~\figurename}
\renewcommand*{\tableformat}{\thetable~\tablename}
\renewcommand*{\captionformat}{~~}
\renewcommand*{\partformat}{\thepart~\partname}
\addto\captionslithuanian{%
  \renewcommand{\figurename}{pav.}%
  \renewcommand{\tablename}{lentelė}%
  \renewcommand{\partname}{dalis}%
  \renewcommand{\contentsname}{Turinys}}
% A sentence that ends on a reference ends on one dot: the
% abbreviation keeps its dot, and \cref / \Cref swallow the full
% stop that follows them ("žr. 9 sk." — never "9 sk..").
\usepackage{letltxmacro}
\makeatletter
\newif\ifltabbr
\newcommand{\ltabbrend}[1]{#1.\global\ltabbrtrue}
% \futurelet, not \@ifnextchar: the space after a reference must stay
\newcommand{\ltgobbledot}{\futurelet\ltnexttoken\ltgobbledotcheck}
\newcommand{\ltgobbledotcheck}{\ifx\ltnexttoken.\expandafter\@gobble\fi}
\AtBeginDocument{%
  \LetLtxMacro\ltorigcref\cref \LetLtxMacro\ltorigCref\Cref
  \DeclareRobustCommand{\cref}[1]{\global\ltabbrfalse\ltorigcref{#1}\ifltabbr\expandafter\ltgobbledot\fi}%
  \DeclareRobustCommand{\Cref}[1]{\global\ltabbrfalse\ltorigCref{#1}\ifltabbr\expandafter\ltgobbledot\fi}}
\makeatother
\newcommand{\ltcref}[2]{%  type, abbreviation without its dot
  \crefformat{#1}{##2##1##3~\ltabbrend{#2}}\Crefformat{#1}{##2##1##3~\ltabbrend{#2}}%
  \crefrangeformat{#1}{##3##1##4--##5##2##6~\ltabbrend{#2}}\Crefrangeformat{#1}{##3##1##4--##5##2##6~\ltabbrend{#2}}%
  \crefmultiformat{#1}{##2##1##3}{ ir~##2##1##3~\ltabbrend{#2}}{, ##2##1##3}{ ir~##2##1##3~\ltabbrend{#2}}%
  \Crefmultiformat{#1}{##2##1##3}{ ir~##2##1##3~\ltabbrend{#2}}{, ##2##1##3}{ ir~##2##1##3~\ltabbrend{#2}}}
\newcommand{\ltcrefword}[2]{%  type, a word in full
  \crefformat{#1}{##2##1##3~#2}\Crefformat{#1}{##2##1##3~#2}%
  \crefrangeformat{#1}{##3##1##4--##5##2##6~#2}\Crefrangeformat{#1}{##3##1##4--##5##2##6~#2}%
  \crefmultiformat{#1}{##2##1##3}{ ir~##2##1##3~#2}{, ##2##1##3}{ ir~##2##1##3~#2}%
  \Crefmultiformat{#1}{##2##1##3}{ ir~##2##1##3~#2}{, ##2##1##3}{ ir~##2##1##3~#2}}
\ltcref{chapter}{sk}
\ltcref{section}{posk}
\ltcref{subsection}{posk}
\ltcref{subsubsection}{posk}
\ltcrefword{appendix}{priedas}
\ltcref{subappendix}{posk}
\ltcref{subsubappendix}{posk}
\ltcref{figure}{pav}
\ltcref{table}{lent}
\ltcref{page}{p}
\newcommand{\crefpairconjunction}{ ir~}
\newcommand{\crefmiddleconjunction}{, }
\newcommand{\creflastconjunction}{ ir~}
\newcommand{\crefpairgroupconjunction}{ ir~}
\newcommand{\crefmiddlegroupconjunction}{, }
\newcommand{\creflastgroupconjunction}{ ir~}
\newcommand{\crefrangeconjunction}{--}
